\documentclass{report}
\usepackage{amsthm}
\usepackage{tcolorbox}
\usepackage{lipsum}
\usepackage{graphicx}
\usepackage{hyperref}
\usepackage{fancyhdr}

% Define theorem environments
\newtheorem{mytheorem}{Theorem}[section]
\newtheorem{mydefinition}[mytheorem]{Definition}

% Define tcolorbox styles
\tcbset{
    myexample/.style={
        enhanced,
        colback=blue!10!white,
        colframe=blue!30!black,
        fonttitle=\bfseries,
        arc=0mm,
        separator sign none,
        leftrule=0pt,
        rightrule=0pt,
        top=4pt,
        bottom=4pt,
        attach boxed title to top left={yshift=-0.1in, xshift=0.15in},
        boxed title style={empty, size=minimal, bottom=1.5pt, top=1.5pt},
        title={Example}
    },
    mytheorem/.style={
        enhanced,
        colback=green!10!white,
        colframe=green!30!black,
        fonttitle=\bfseries,
        arc=0mm,
        separator sign none,
        leftrule=0pt,
        rightrule=0pt,
        top=4pt,
        bottom=4pt,
        attach boxed title to top left={yshift=-0.1in, xshift=0.15in},
        boxed title style={empty, size=minimal, bottom=1.5pt, top=1.5pt},
        title={Theorem}
    },
    mydefinition/.style={
        enhanced,
        colback=orange!10!white,
        colframe=orange!30!black,
        fonttitle=\bfseries,
        arc=0mm,
        separator sign none,
        leftrule=0pt,
        rightrule=0pt,
        top=4pt,
        bottom=4pt,
        attach boxed title to top left={yshift=-0.1in, xshift=0.15in},
        boxed title style={empty, size=minimal, bottom=1.5pt, top=1.5pt},
        title={Definition}
    },
    mynote/.style={
        enhanced,
        colback=gray!10!white,
        colframe=gray!30!black,
        fonttitle=\bfseries,
        arc=0mm,
        separator sign none,
        leftrule=0pt,
        rightrule=0pt,
        top=4pt,
        bottom=4pt,
        attach boxed title to top left={yshift=-0.1in, xshift=0.15in},
        boxed title style={empty, size=minimal, bottom=1.5pt, top=1.5pt},
        title={Note}
    }
}

% Define tcolorbox environments
\newtcolorbox{myexample}[1][]{
    myexample,
    #1
}

\newtcolorbox{mytheorem}[1][]{
    mytheorem,
    #1
}

\newtcolorbox{mydefinition}[1][]{
    mydefinition,
    #1
}

\newtcolorbox{mynote}[1][]{
    mynote,
    #1
}

% Header and footer settings
\pagestyle{fancy}
\fancyhf{}
\fancyhead[L]{Tannon Warnick}
\fancyhead[C]{Class Name}
\fancyhead[R]{\thepage}
\renewcommand{\headrulewidth}{1pt} % Line under the header
\renewcommand{\footrulewidth}{0pt} % No line in the footer

\begin{document}

\begin{titlepage}
    \centering
    \includegraphics[width=0.5\textwidth]{University_logo.jpeg}\par\vspace{1cm}
    {\scshape\Large Linear Algebra Notes \par}
    \vspace{1.5cm}
    {\huge\bfseries Tannon Warnick \par}
    \vspace{2cm}
    {\Large Class Name\par}
    \vfill
    {\large \today\par}
\end{titlepage}

\tableofcontents
\newpage

\chapter{Introduction}
\section{Overview}
\lipsum[1-2]

\section{Motivation}
\lipsum[3-4]

\chapter{Vectors and Spaces}
\section{Vector Definitions}
\begin{mydefinition}
A vector is an element of a vector space.
\end{mydefinition}

\section{Vector Spaces}
\begin{mytheorem}
Every vector space has a basis.
\end{mytheorem}

\begin{myexample}
The set of all $n \times n$ matrices forms a vector space.
\tcblower
sadfsdfawe
\end{myexample}

\chapter{Linear Transformations}
\section{Definition}
\begin{mydefinition}
A linear transformation is a mapping between vector spaces that preserves vector addition and scalar multiplication.
\end{mydefinition}

\section{Properties}
\begin{mytheorem}
A linear transformation is injective if and only if its kernel is trivial.
\end{mytheorem}

\chapter{Conclusion}
\lipsum[5-6]

\end{document}
